% ============================================================
% ECON 320 — Introduction to Mathematical Economics
% Quiz 1 — INSTRUCTOR ANSWER KEY
% Fall 2026  |  Muhebullah Karimzada  |  UNBC
% ============================================================
% DO NOT DISTRIBUTE TO STUDENTS
% ============================================================
\documentclass[11pt]{article}
\usepackage[margin=1in,top=1.2in,bottom=1.2in]{geometry}
\usepackage[utf8]{inputenc}
\usepackage[T1]{fontenc}
\usepackage{lmodern}
\usepackage{amsmath,amssymb,bm}
\usepackage{booktabs,array,multirow,adjustbox}
\usepackage{enumitem}
\usepackage{xcolor}
\usepackage{fancyhdr}
\setlength{\headheight}{22pt}
\addtolength{\topmargin}{-10pt}
\usepackage{tcolorbox}
\usepackage{titlesec}
\usepackage{hyperref}
\tcbuselibrary{skins}
\hypersetup{colorlinks=false,hidelinks}

%% ── Colours ────────────────────────────────────────────────
\definecolor{UNBCDark}   {RGB}{0,74,37}
\definecolor{UNBCGold}   {RGB}{189,155,96}
\definecolor{SoftGray}   {RGB}{244,246,248}
\definecolor{CorrectGreen}{RGB}{0,120,60}
\definecolor{AccentRed}  {RGB}{192,57,43}
\definecolor{AccentBlue} {RGB}{41,98,255}
\definecolor{AccentOrange}{RGB}{230,126,34}

%% ── Header / Footer ────────────────────────────────────────
\pagestyle{fancy}
\fancyhf{}
\renewcommand{\headrulewidth}{0.5pt}
\renewcommand{\headrule}{\hbox to\headwidth{\color{AccentRed}%
  \leaders\hrule height\headrulewidth\hfill}}
\lhead{\small\textbf{\color{UNBCDark}ECON 320 --- Quiz 1}%
       \;\textcolor{UNBCGold}{$|$}\;%
       \textcolor{AccentRed}{\bfseries INSTRUCTOR ANSWER KEY}}
\rhead{\small\color{UNBCDark}Fall 2026}
\cfoot{\small\thepage}

%% ── Section style ──────────────────────────────────────────
\titleformat{\section}[block]
  {\large\bfseries\color{UNBCDark}}{}
  {0em}{}
  [{\color{UNBCGold}\titlerule[0.8pt]}]
\titlespacing*{\section}{0pt}{1.2em}{0.6em}

\titleformat{\subsection}[runin]
  {\bfseries\color{UNBCDark}}{}
  {0em}{}[.\enspace]
\titlespacing*{\subsection}{0pt}{0.8em}{0.3em}

%% ── Custom boxes ───────────────────────────────────────────
\newtcolorbox{confidential}{
  colback=AccentRed!8, colframe=AccentRed,
  left=4mm, right=4mm, top=3mm, bottom=3mm,
  boxrule=0.8pt, arc=3mm}

\newtcolorbox{correctbox}{
  colback=CorrectGreen!10, colframe=CorrectGreen,
  left=3mm, right=3mm, top=1.5mm, bottom=1.5mm,
  boxrule=0.6pt, arc=2mm}

\newtcolorbox{solutionbox}{
  colback=AccentBlue!5, colframe=AccentBlue!60,
  title={\small\bfseries\color{UNBCDark}Solution},
  left=4mm, right=4mm, top=2mm, bottom=2mm,
  boxrule=0.6pt, arc=2mm}

\newtcolorbox{rubricbox}{
  colback=AccentOrange!8, colframe=AccentOrange,
  title={\small\bfseries Grading Rubric},
  left=3mm, right=3mm, top=2mm, bottom=2mm,
  boxrule=0.5pt, arc=2mm}

\newtcolorbox{databox}{
  colback=SoftGray, colframe=UNBCDark!40,
  left=4mm, right=4mm, top=3mm, bottom=3mm,
  boxrule=0.4pt, arc=2mm}

% ============================================================
\begin{document}
% ============================================================



\vspace{0.5em}

\begin{center}
{\LARGE\bfseries\color{UNBCDark}ECON 320 \;--\; Introduction to Mathematical Economics}\\[0.5em]
{\large\bfseries Quiz 1\quad\textcolor{UNBCGold}{$|$}\quad Fall 2026}\\[0.2em]
{\small University of Northern British Columbia\quad Instructor: M.\,Karimzada}
\end{center}

\vspace{0.3em}
{\color{UNBCGold}\hrule height 1.5pt}
\vspace{0.5em}



\vspace{0.6em}

%% ════════════════════════════════════════════════════════════
\section{Part A: Multiple Choice --- Answer Key}
%% ════════════════════════════════════════════════════════════



\bigskip

%% ── Q1 ──────────────────────────────────────────────────────
\noindent\textbf{Q1.}\quad\textbf{Correct Answer: B}

\begin{correctbox}
\textbf{B)}\; $\{1,\,3,\,4,\,5,\,7,\,8\}$
\end{correctbox}

\medskip
\noindent\textbf{Explanation.}\enspace
First compute $A \cap B = \{1,2,5,6\} \cap \{2,3,6,7\} = \{2,6\}$.
Then $(A \cap B)^c = U \setminus \{2,6\} = \{1,3,4,5,7,8\}$.
Option A is just $A \cap B$ itself. Option C is $A \cup B$. Option D is neither.

\bigskip

%% ── Q2 ──────────────────────────────────────────────────────
\noindent\textbf{Q2.}\quad\textbf{Correct Answer: C}

\begin{correctbox}
\textbf{C)}\; $32$
\end{correctbox}

\medskip
\noindent\textbf{Explanation.}\enspace
The power set of a set with $n$ elements has $2^n$ elements.
Here $n = 5$, so $|\mathcal{P}(A)| = 2^5 = 32$.
Option A ($10 = \binom{5}{2}$) counts only 2-element subsets.
Option B ($25 = 5^2$) confuses exponent and base.
Option D ($64 = 2^6$) uses the wrong exponent.

\bigskip

%% ── Q3 ──────────────────────────────────────────────────────
\noindent\textbf{Q3.}\quad\textbf{Correct Answer: B}

\begin{correctbox}
\textbf{B)}\; $f(x) = e^x$
\end{correctbox}

\medskip
\noindent\textbf{Explanation.}\enspace
$f(x) = e^x$ is strictly increasing on $\mathbb{R}$, so it is injective.
Its range is $(0,\infty) \neq \mathbb{R}$, so it is not surjective onto $\mathbb{R}$.
\begin{itemize}[nosep,leftmargin=1.3em]
  \item \textbf{A)} $x^2$: not injective since $f(-1)=f(1)=1$.
  \item \textbf{C)} $x^3$: bijective (both injective and surjective onto $\mathbb{R}$).
  \item \textbf{D)} $\sin x$: neither injective ($\sin 0 = \sin \pi$) nor surjective onto $\mathbb{R}$ (range $= [-1,1]$).
\end{itemize}

\bigskip

%% ── Q4 ──────────────────────────────────────────────────────
\noindent\textbf{Q4.}\quad\textbf{Correct Answer: A}

\begin{correctbox}
\textbf{A)}\; $\displaystyle f^{-1}(x) = \frac{x+3}{x-2}$
\end{correctbox}

\medskip
\noindent\textbf{Explanation.}\enspace
Set $y = \dfrac{2x+3}{x-1}$ and solve for $x$:
\[
  y(x-1) = 2x+3
  \;\Rightarrow\; xy - y = 2x+3
  \;\Rightarrow\; xy - 2x = y+3
  \;\Rightarrow\; x(y-2) = y+3
  \;\Rightarrow\; x = \frac{y+3}{y-2}.
\]
So $f^{-1}(x) = \dfrac{x+3}{x-2}$, defined for $x \neq 2$.
The other options result from algebraic errors in rearranging the equation.

\bigskip

%% ── Q5 ──────────────────────────────────────────────────────
\noindent\textbf{Q5.}\quad\textbf{Correct Answer: B}

\begin{correctbox}
\textbf{B)}\; $\sqrt{10}$
\end{correctbox}

\medskip
\noindent\textbf{Explanation.}\enspace
Apply $f$ first, then $g$:
\[
  (g \circ f)(3) = g(f(3)) = g(3^2 - 1) = g(8) = \sqrt{8+2} = \sqrt{10}.
\]
Option A ($2\sqrt{2}$) would arise from computing $g(f(\sqrt{3}))$. Options C and D result from order-of-composition errors.



%% ════════════════════════════════════════════════════════════
\section{Part B: Numerical Problems --- Full Solution}
%% ════════════════════════════════════════════════════════════

%% ── Q6 ──────────────────────────────────────────────────────
\noindent\textbf{Q6.}\quad\textbf{Set Operations.}\quad(5 marks)

\medskip
\begin{databox}
\small
$U = \{1,2,3,4,5,6,7,8,9,10\}$,\quad
$A = \{1,3,5,7,9\}$,\quad
$B = \{1,2,3,4,5\}$,\quad
$C = \{3,5,6,7,8\}$.
\end{databox}



\begin{solutionbox}
\textbf{(a)} $A \cap B$: elements in both $A$ and $B$:
\[A \cap B = \{1,3,5\}.\]
$A \setminus B$: elements in $A$ but not in $B$:
\[A \setminus B = \{7,9\}.\]

\medskip
\textbf{(b)} First, $A \cup B = \{1,2,3,4,5,7,9\}$.
\[(A \cup B)^c = U \setminus \{1,2,3,4,5,7,9\} = \{6,8,10\}.\]

\medskip
\textbf{(c)} Compute $A^c$ and $B^c$ separately, then intersect:
\[A^c = \{2,4,6,8,10\}, \qquad B^c = \{6,7,8,9,10\}.\]
\[A^c \cap B^c = \{6,8,10\}.\]
Since $(A \cup B)^c = \{6,8,10\} = A^c \cap B^c$, De Morgan's Law is verified. $\checkmark$

\medskip
\textbf{(d)} From (a), $A \cap B = \{1,3,5\}$. Now intersect with $C = \{3,5,6,7,8\}$:
\[A \cap B \cap C = \{3,5\}.\]
Since $|A \cap B \cap C| = 2$:
\[\bigl|\mathcal{P}(A \cap B \cap C)\bigr| = 2^2 = \mathbf{4}.\]
\end{solutionbox}



\bigskip

%% ── Q7 ──────────────────────────────────────────────────────
\noindent\textbf{Q7.}\quad\textbf{Inverse Functions.}\quad(5 marks)

\medskip
\begin{databox}
\small $\displaystyle f(x) = \frac{2x+3}{x-1}$,\quad $x \neq 1$.
\end{databox}

\medskip

\begin{solutionbox}
\textbf{(a) Domain and Range.}

\textit{Domain:} $f$ is undefined at $x = 1$, so $\operatorname{dom}(f) = \mathbb{R} \setminus \{1\}$.

\textit{Range:} Set $y = \dfrac{2x+3}{x-1}$ and solve for $x$:
$y(x-1) = 2x+3 \;\Rightarrow\; x(y-2) = y+3 \;\Rightarrow\; x = \dfrac{y+3}{y-2}$.

This requires $y \neq 2$, so $\operatorname{ran}(f) = \mathbb{R} \setminus \{2\}$.

\medskip
\textbf{(b) Inverse Function.}

From part (a):
\[
  f^{-1}(x) = \frac{x+3}{x-2}, \qquad \operatorname{dom}(f^{-1}) = \mathbb{R} \setminus \{2\}.
\]

\medskip
\textbf{(c) Verification: $f(f^{-1}(x)) = x$.}
\[
  f\!\left(\frac{x+3}{x-2}\right)
  = \frac{2\cdot\dfrac{x+3}{x-2}+3}{\dfrac{x+3}{x-2}-1}
  = \frac{\dfrac{2(x+3)+3(x-2)}{x-2}}{\dfrac{(x+3)-(x-2)}{x-2}}
  = \frac{2x+6+3x-6}{x+3-x+2}
  = \frac{5x}{5} = x. \quad\checkmark
\]
\end{solutionbox}



\newpage

%% ── Q8 ──────────────────────────────────────────────────────
\noindent\textbf{Q8.}\quad\textbf{Profit Function Analysis.}\quad(5 marks)

\medskip
\begin{databox}
\small $\pi(x) = -2x^2 + 12x - 10$,\quad $x \geq 0$.
\end{databox}

\medskip

\begin{solutionbox}
\textbf{(a) Break-even quantities: $\pi(x) = 0$.}
\[
  -2x^2 + 12x - 10 = 0
  \;\Rightarrow\; x^2 - 6x + 5 = 0
  \;\Rightarrow\; (x-1)(x-5) = 0.
\]
Break-even outputs: $\mathbf{x = 1}$ and $\mathbf{x = 5}$.

\medskip
\textbf{(b) Profit-maximising quantity and maximum profit.}
\[
  \pi'(x) = -4x + 12 = 0 \;\Rightarrow\; x^* = 3.
\]
\[
  \pi^* = \pi(3) = -2(9) + 12(3) - 10 = -18 + 36 - 10 = \mathbf{8}.
\]

\medskip
\textbf{(c) Concavity.}
\[
  \pi''(x) = -4 < 0 \quad\text{for all } x \geq 0.
\]
Since $\pi''(x) < 0$ everywhere, $\pi$ is \textbf{strictly concave}.
This confirms $x^* = 3$ is a global maximum and reflects diminishing marginal profit.

\medskip
\textbf{(d) Range of $\pi$ on $[0,\infty)$.}

Complete the square: $\pi(x) = -2(x-3)^2 + 8$.

At $x = 0$: $\pi(0) = -10$. Maximum at $x^* = 3$: $\pi(3) = 8$. As $x \to \infty$: $\pi(x) \to -\infty$.

\[\text{Range of } \pi = (-\infty,\; 8].\]
\end{solutionbox}



\vspace{1em}
{\color{UNBCGold}\hrule height 0.6pt}



\vspace{1em}
{\color{UNBCGold}\hrule height 0.6pt}
\begin{center}
\small\itshape
--- End of Answer Key ---\quad ECON 320, Fall 2026, Quiz 1
\end{center}

\end{document}
